\documentclass[11pt,a4paper]{article}
\usepackage[T1]{fontenc}
\usepackage[utf8]{inputenc}
\usepackage[french]{babel}
\usepackage{lmodern}
\usepackage{microtype}
\usepackage[margin=2.2cm]{geometry}
\usepackage{amsmath,amssymb,amsthm,mathtools}
\usepackage{icomma}
\usepackage{booktabs,tabularx,array}
\usepackage{xcolor}
\usepackage{tikz}
\usetikzlibrary{patterns,positioning}
\usepackage{pgfplots}
\pgfplotsset{compat=1.18}
\usepackage[most]{tcolorbox}
\usepackage{enumitem}
\usepackage[hidelinks]{hyperref}
\usepackage{caption}
\usepackage{float}
\usepackage{pdflscape}
\usepackage{pifont}

\definecolor{nyt}{HTML}{1F3B63}
\definecolor{wp}{HTML}{17695B}
\definecolor{gov}{HTML}{B42318}
\definecolor{nat}{HTML}{9A6212}
\definecolor{regnone}{HTML}{C6C2B8}
\definecolor{regnosuit}{HTML}{CBE2C3}
\definecolor{regnyt}{HTML}{F0C5BD}
\definecolor{regboth}{HTML}{F1D69B}

\newtcolorbox{cours}{colback=nat!6,colframe=nat!70,boxrule=0.5pt,arc=2pt,left=6pt,right=6pt,top=4pt,bottom=4pt,fonttitle=\bfseries\small,title=Lien avec le cours}
\newtcolorbox{resultat}{colback=wp!6,colframe=wp!70,boxrule=0.5pt,arc=2pt,left=6pt,right=6pt,top=4pt,bottom=4pt}
\newtcolorbox{intuition}{colback=nyt!4,colframe=nyt!50,boxrule=0.4pt,arc=2pt,left=6pt,right=6pt,top=3pt,bottom=3pt,fonttitle=\bfseries\small,title=Intuition}

\theoremstyle{definition}
\newtheorem{lemme}{Lemme}
\newtheorem{proposition}{Proposition}
\newcommand{\ind}[1]{\mathbf{1}_{\{\text{#1}\}}}
\newcommand{\E}{\mathbb{E}}
\setlength{\parindent}{0pt}
\setlength{\parskip}{5pt}
\captionsetup{font=small,labelfont=bf}

\newtcolorbox{erreur}{colback=gov!5,colframe=gov!70,boxrule=0.5pt,arc=2pt,left=6pt,right=6pt,top=4pt,bottom=4pt,fonttitle=\bfseries\small,title=Erreur}
\newtcolorbox{discuter}{colback=nat!6,colframe=nat!70,boxrule=0.5pt,arc=2pt,left=6pt,right=6pt,top=4pt,bottom=4pt,fonttitle=\bfseries\small,title=À discuter}
\newtcolorbox{simple}{colback=nyt!4,colframe=nyt!50,boxrule=0.4pt,arc=2pt,left=6pt,right=6pt,top=3pt,bottom=3pt,fonttitle=\bfseries\small,title=En mots}
\newcommand{\T}[1]{\textcolor{nyt}{#1}}
\newcommand{\W}[1]{\textcolor{wp}{#1}}
\newcommand{\N}[1]{\textcolor{gov}{#1}}

\title{\textbf{Pentagon Papers}\\[3pt]\large Relecture de l'arbre du collègue : erreurs, corrections et calculs}
\author{Travail de théorie des jeux, groupe 5--6 · ECGEB366}
\date{}

\begin{document}
\maketitle
\vspace{-1.2cm}

\begin{resultat}
\textbf{En bref.}
\begin{itemize}[leftmargin=*,itemsep=1pt,topsep=2pt]
\item L'arbre du collègue est \textbf{le bon} : mêmes joueurs, même ordre, mêmes sept feuilles que notre modèle.
\item Quatre feuilles sont justes ($O_3$, $O_4$, $O_5$, $O_6$). Les erreurs sont dans $O_1$ (trois), $O_2$ (une) et $O_7$ (une).
\item Avec ses valeurs corrigées et sans la part déjà publiée $\theta$, \textbf{rien n'est publié} ($O_7$).
\item Avec $\theta=0,3$, \textbf{on retrouve 1971} ($O_1$).
\end{itemize}
\end{resultat}

\begin{discuter}
\textbf{Attention aux lettres.} Dans ce document, on garde les lettres du collègue : $p$ = chance que le \textbf{juge local} refuse la censure ; $q$ = chance que la \textbf{Cour suprême} donne raison aux journaux. C'est l'inverse de nos autres PDF. Le collègue n'a pas fixé de valeurs : les exemples chiffrés utilisent $p=0,3$ et $q=0,65$.
\end{discuter}

\section{Les valeurs}

\subsection{Ce que le dessin utilise, et ce que la liste dit}
En lisant les feuilles du dessin, on retrouve les valeurs employées. Elles ne correspondent pas toutes à la liste annoncée.
\begin{center}
\footnotesize
\begin{tabular}{@{}llll@{}}
\toprule
Paramètre & Symbole & Dans le dessin & Dans la liste \\
\midrule
Bénéfice (les deux journaux) & $B$ & $10$ & $10$ \\
Coût du procès, Times & $c_T$ & $1$ \quad ($9=10-1$) & $1$ \\
Coût du procès, Post & $c_W$ & $\mathbf{3}$ \quad ($7=10-3$ ; $-8=-3-5$) & $\mathbf{5}$ \\
Sanction si censuré & $S$ & $\mathbf{5}$ \quad ($-6=-1-5$) & non mentionnée \\
Dommage de Nixon si tout sort & $D$ & $\mathbf{5}$ \quad (feuilles à $-5$ et $-6$) & $\mathbf{200}$ \\
Coût d'un procès pour Nixon & $k$ & $1$ & $1$ \\
\bottomrule
\end{tabular}
\end{center}

\subsection{Les valeurs retenues ici}
On suit la liste, et on garde la sanction du dessin :
\[
B_T=B_W=10,\quad c_T=1,\quad c_W=5,\quad S_T=S_W=5,\quad D=200,\quad k_T=k_W=1 .
\]
\begin{simple}
On écrit tous les paramètres en \textbf{positif} et on met les signes « $-$ » dans les équations. Le dessin écrit « $\cdot-1$ » ou « $\cdot-6$ » : cela marche, mais les doubles moins font facilement des erreurs.
\end{simple}

\section{Les erreurs, feuille par feuille}
\begin{landscape}
\begin{center}{\large\bfseries L'arbre du collègue, annoté}\\[2pt]{\small Gains dans l'ordre \textcolor{nyt}{\textbf{Times}}, \textcolor{wp}{\textbf{Post}}, \textcolor{gov}{\textbf{Nixon}}. $p$ = juge local, $q$ = Cour suprême. \textcolor{wp}{\ding{51}} juste \quad \textcolor{gov}{\ding{55}} erreur \quad \textcolor{nat}{?} à discuter}\end{center}\vspace{-4pt}
\noindent\resizebox{\linewidth}{!}{\begin{tikzpicture}[x=1cm,y=1cm,lab/.style={font=\scriptsize,fill=white,inner sep=1pt,sloped},eq/.style={line width=1.8pt,draw=gov},br/.style={line width=0.6pt,draw=black!55}]
\draw[br] (0.4,-7.928125)--(2.6,-4.56875) node[lab,midway,above]{Publier};
\draw[br] (0.4,-7.928125)--(5.0,-11.2875) node[lab,midway,below]{Ne pas publier};
\draw[br] (2.6,-4.56875)--(5.0,-2.6875) node[lab,midway,above]{Poursuivre};
\draw[br] (2.6,-4.56875)--(9.6,-6.449999999999999) node[lab,midway,below]{Laisser faire};
\draw[br] (5.0,-2.6875)--(7.4,-1.075) node[lab,midway,above]{Publier};
\draw[br] (5.0,-2.6875)--(9.6,-4.3) node[lab,midway,below]{Ne pas publier};
\draw[br] (7.4,-1.075)--(9.6,0.0) node[lab,midway,above]{Poursuivre};
\draw[br] (7.4,-1.075)--(9.6,-2.15) node[lab,midway,below]{Laisser faire};
\draw[br] (5.0,-11.2875)--(7.4,-9.675) node[lab,midway,above]{Publier};
\draw[br] (5.0,-11.2875)--(9.6,-12.899999999999999) node[lab,midway,below]{Ne pas publier};
\draw[br] (7.4,-9.675)--(9.6,-8.6) node[lab,midway,above]{Poursuivre};
\draw[br] (7.4,-9.675)--(9.6,-10.75) node[lab,midway,below]{Laisser faire};
\fill[nyt] (0.4,-7.928125) circle (4pt); \node[font=\small\bfseries,text=nyt,above left=3pt and -2pt] at (0.4,-7.928125) {Times};
\fill[gov] (2.6,-4.56875) circle (4pt); \node[font=\small\bfseries,text=gov,above left=3pt and -2pt] at (2.6,-4.56875) {Nixon};
\fill[wp] (5.0,-2.6875) circle (4pt); \node[font=\small\bfseries,text=wp,above left=3pt and -2pt] at (5.0,-2.6875) {Post};
\fill[gov] (7.4,-1.075) circle (4pt); \node[font=\small\bfseries,text=gov,above left=3pt and -2pt] at (7.4,-1.075) {Nixon};
\fill[wp] (5.0,-11.2875) circle (4pt); \node[font=\small\bfseries,text=wp,above left=3pt and -2pt] at (5.0,-11.2875) {Post};
\fill[gov] (7.4,-9.675) circle (4pt); \node[font=\small\bfseries,text=gov,above left=3pt and -2pt] at (7.4,-9.675) {Nixon};
\fill (9.6,0.0) circle (2pt);
\node[anchor=west,draw=black!35,fill=white,rounded corners=3pt,text width=15.2cm,align=left,font=\footnotesize,inner sep=4pt] at (9.85,0.0) {\textbf{$O_1$ · les deux journaux poursuivis}\\\textcolor{nyt}{\textbf{Times} : $p\cdot9+(1-p)\cdot(-6)$\ \textcolor{gov}{\ding{55}\ \textit{oublie la Cour suprême : manque $(1-p)q\cdot9$}}}\\\textcolor{wp}{\textbf{Post} : $pq\cdot7+(1-p)(1-q)\cdot(-8)$\ \textcolor{gov}{\ding{55}\ \textit{les chances ne font que 44\,\%}}}\\\textcolor{gov}{\textbf{Nixon} : $(1-p)(1-q)\cdot(-1)+pq\cdot(-6)$\ \textcolor{gov}{\ding{55}\ \textit{chances fausses et 2 procès : $-2$ et $-202$}}}};
\fill (9.6,-2.15) circle (2pt);
\node[anchor=west,draw=black!35,fill=white,rounded corners=3pt,text width=15.2cm,align=left,font=\footnotesize,inner sep=4pt] at (9.85,-2.15) {\textbf{$O_2$ · Times poursuivi, Post libre}\\\textcolor{nyt}{\textbf{Times} : $p\cdot9+(1-p)\cdot(-6)$\ \textcolor{nat}{?\ \textit{le Post a tout publié : 9 ?}}}\\\textcolor{wp}{\textbf{Post} : $10$\ \textcolor{nat}{?\ \textit{0 si le juge libère le Times : $(1-p)\cdot10$}}}\\\textcolor{gov}{\textbf{Nixon} : $(1-p)\cdot(-1)+p\cdot(-6)$\ \textcolor{gov}{\ding{55}\ \textit{tout sort dans les deux cas : $-201$}}}};
\fill (9.6,-4.3) circle (2pt);
\node[anchor=west,draw=black!35,fill=white,rounded corners=3pt,text width=15.2cm,align=left,font=\footnotesize,inner sep=4pt] at (9.85,-4.3) {\textbf{$O_3$ · Times poursuivi, Post muet}\\\textcolor{nyt}{\textbf{Times} : $p\cdot9+(1-p)\cdot(-6)$\ \textcolor{wp}{\ding{51}}}\\\textcolor{wp}{\textbf{Post} : $0$\ \textcolor{wp}{\ding{51}}}\\\textcolor{gov}{\textbf{Nixon} : $(1-p)\cdot(-1)+p\cdot(-6)$\ \textcolor{wp}{\ding{51}}}};
\fill (9.6,-6.449999999999999) circle (2pt);
\node[anchor=west,draw=black!35,fill=white,rounded corners=3pt,text width=15.2cm,align=left,font=\footnotesize,inner sep=4pt] at (9.85,-6.449999999999999) {\textbf{$O_4$ · Times publie, pas de procès}\\\textcolor{nyt}{\textbf{Times} : $10$\ \textcolor{wp}{\ding{51}}}\\\textcolor{wp}{\textbf{Post} : $0$\ \textcolor{wp}{\ding{51}}}\\\textcolor{gov}{\textbf{Nixon} : $-5$\ \textcolor{nat}{?\ \textit{avec $D=200$ : $-200$}}}};
\fill (9.6,-8.6) circle (2pt);
\node[anchor=west,draw=black!35,fill=white,rounded corners=3pt,text width=15.2cm,align=left,font=\footnotesize,inner sep=4pt] at (9.85,-8.6) {\textbf{$O_5$ · Post seul, poursuivi}\\\textcolor{nyt}{\textbf{Times} : $0$\ \textcolor{wp}{\ding{51}}}\\\textcolor{wp}{\textbf{Post} : $p\cdot7+(1-p)\cdot(-8)$\ \textcolor{nat}{?\ \textit{juste avec un coût de 3 ; avec 5 : 5 et $-10$}}}\\\textcolor{gov}{\textbf{Nixon} : $(1-p)\cdot(-1)+p\cdot(-6)$\ \textcolor{wp}{\ding{51}}}};
\fill (9.6,-10.75) circle (2pt);
\node[anchor=west,draw=black!35,fill=white,rounded corners=3pt,text width=15.2cm,align=left,font=\footnotesize,inner sep=4pt] at (9.85,-10.75) {\textbf{$O_6$ · Post seul, pas de procès}\\\textcolor{nyt}{\textbf{Times} : $0$\ \textcolor{wp}{\ding{51}}}\\\textcolor{wp}{\textbf{Post} : $10$\ \textcolor{wp}{\ding{51}}}\\\textcolor{gov}{\textbf{Nixon} : $-5$\ \textcolor{nat}{?\ \textit{avec $D=200$ : $-200$}}}};
\fill (9.6,-12.899999999999999) circle (2pt);
\node[anchor=west,draw=black!35,fill=white,rounded corners=3pt,text width=15.2cm,align=left,font=\footnotesize,inner sep=4pt] at (9.85,-12.899999999999999) {\textbf{$O_7$ · personne ne publie}\\\textcolor{nyt}{\textbf{Times} : $0$\ \textcolor{wp}{\ding{51}}}\\\textcolor{wp}{\textbf{Post} : $0$\ \textcolor{wp}{\ding{51}}}\\\textcolor{gov}{\textbf{Nixon} : $10$\ \textcolor{nat}{?\ \textit{meilleur cas : 0 suffit ; si +10, l\'ajouter aussi aux censures (9 et 8)}}}};
\end{tikzpicture}}
\end{landscape}


\subsection{Le tableau de contrôle}
\begin{center}
\small
\begin{tabularx}{\textwidth}{@{}l>{\raggedright\arraybackslash}Xl@{}}
\toprule
Feuille & Ce que le collègue a écrit (Times ; Post ; Nixon) & Verdict \\
\midrule
$O_4$ Times publie, pas de procès & $(10\ ;\ 0\ ;\ -5)$ & juste \\
$O_3$ Times poursuivi, Post muet & $p\cdot9+(1-p)(-6)$ ; $0$ ; $(1-p)(-1)+p(-6)$ & juste \\
$O_5$ Post seul, poursuivi & $0$ ; $p\cdot7+(1-p)(-8)$ ; $(1-p)(-1)+p(-6)$ & juste (avec $c_W=3$) \\
$O_6$ Post seul, pas de procès & $(0\ ;\ 10\ ;\ -5)$ & juste \\
$O_7$ Personne ne publie & $(0\ ;\ 0\ ;\ \mathbf{10})$ & \textbf{incohérent} \\
$O_2$ Times poursuivi, Post libre & $p\cdot9+(1-p)(-6)$ ; $10$ ; $(1-p)(-1)+p(-6)$ & \textbf{1 erreur} \\
$O_1$ Les deux poursuivis & $p\cdot9+(1-p)(-6)$ ; $pq\cdot7+(1-p)(1-q)(-8)$ ; $(1-p)(1-q)(-1)+pq(-6)$ & \textbf{3 erreurs} \\
\bottomrule
\end{tabularx}
\end{center}

\subsection{$O_7$ : Nixon ne peut pas avoir $+10$}
\begin{erreur}
Partout ailleurs, le meilleur cas de Nixon est « rien ne sort, aucun procès », qui vaut $0$ avec son équation $-D\sigma-k$. Un $+10$ ici n'a pas de justification.
\end{erreur}
Conséquence : aucune sur l'équilibre, car Nixon ne décide jamais dans cette branche. Mais l'incohérence se verrait dans le rapport. \textbf{Correction} : $(0\ ;\ 0\ ;\ 0)$.

\subsection{$O_2$ : Nixon oublie que le Post a tout publié}
\begin{erreur}
Le collègue écrit pour Nixon $(1-p)(-1)+p(-6)$ : quand le Times est censuré ($1-p$), Nixon n'aurait que le coût du procès. Or dans $O_2$ le Post publie librement : \textbf{tout sort, quel que soit le verdict}. Nixon a donc toujours $-D-k_T=-201$.
\end{erreur}
\begin{discuter}
\textbf{Le Times dans $O_2$.} Le collègue lui donne la même chose qu'en $O_3$. Mais si le Post a tout publié, la censure du Times ne sert plus à rien : dans notre modèle, le Times a $B_T-c_T=9$ dans tous les cas. On retient ce choix ici.
\end{discuter}

\subsection{$O_1$ : trois erreurs}
\begin{erreur}
\textbf{1. Le Times oublie la Cour suprême.} Les deux affaires vont \emph{ensemble} devant la Cour suprême : le Times peut donc aussi y gagner. Sa formule doit contenir $q$ :
\[
p\cdot9+(1-p)\,q\cdot9+(1-p)(1-q)\cdot(-6).
\]
\end{erreur}
\begin{erreur}
\textbf{2. Les chances du Post ne font pas 100\,\%.} Le collègue n'utilise que deux cas, $pq$ et $(1-p)(1-q)$. Avec $p=0,3$ et $q=0,65$ : $0,195+0,245=0,44$. Il manque $56\,\%$ des cas.

Le terme $pq$ voudrait dire « gagner au juge local \emph{et} à la Cour suprême ». Mais si le juge local refuse la censure, l'affaire s'arrête : la Cour suprême n'est jamais saisie. Les trois chemins corrects sont :
\[
\underbrace{p}_{\text{juge local refuse}}\ +\ \underbrace{(1-p)\,q}_{\text{censure, puis Cour pour les journaux}}\ +\ \underbrace{(1-p)(1-q)}_{\text{censure confirmée}}\ =\ 1 .
\]
\end{erreur}
\begin{erreur}
\textbf{3. Nixon : mêmes chances fausses, et un procès oublié.} Dans $O_1$, Nixon poursuit \emph{deux} journaux. Il faut $-2$ (pas $-1$) si la censure tient et $-202$ (pas $-6$ ou $-201$) si tout sort.
\end{erreur}
\begin{discuter}
\textbf{Le Post si le juge local libère le Times.} Le collègue lui donne $7$. Mais dans ce cas le Times publie tout lui-même : le relais du Post ne sert plus et il n'est pas jugé. Notre modèle lui donne $0$. On retient ce choix ici.
\end{discuter}

\section{Les équations corrigées (sans \texorpdfstring{$\theta$}{theta})}
\begin{landscape}
\begin{center}{\large\bfseries L'arbre corrigé, en lettres}\\[2pt]{\small Gains dans l'ordre \textcolor{nyt}{\textbf{Times}}, \textcolor{wp}{\textbf{Post}}, \textcolor{gov}{\textbf{Nixon}}. $p$ = juge local, $q$ = Cour suprême. Version générale avec $\theta$ ; sans $\theta$, on pose $\theta=0$ et les termes $\theta B$ et $\theta D$ disparaissent.}\end{center}\vspace{-4pt}
\noindent\resizebox{\linewidth}{!}{\begin{tikzpicture}[x=1cm,y=1cm,lab/.style={font=\scriptsize,fill=white,inner sep=1pt,sloped},eq/.style={line width=1.8pt,draw=gov},br/.style={line width=0.6pt,draw=black!55}]
\draw[br] (0.4,-7.928125)--(2.6,-4.56875) node[lab,midway,above]{Publier};
\draw[br] (0.4,-7.928125)--(5.0,-11.2875) node[lab,midway,below]{Ne pas publier};
\draw[br] (2.6,-4.56875)--(5.0,-2.6875) node[lab,midway,above]{Poursuivre};
\draw[br] (2.6,-4.56875)--(9.6,-6.449999999999999) node[lab,midway,below]{Laisser faire};
\draw[br] (5.0,-2.6875)--(7.4,-1.075) node[lab,midway,above]{Publier};
\draw[br] (5.0,-2.6875)--(9.6,-4.3) node[lab,midway,below]{Ne pas publier};
\draw[br] (7.4,-1.075)--(9.6,0.0) node[lab,midway,above]{Poursuivre};
\draw[br] (7.4,-1.075)--(9.6,-2.15) node[lab,midway,below]{Laisser faire};
\draw[br] (5.0,-11.2875)--(7.4,-9.675) node[lab,midway,above]{Publier};
\draw[br] (5.0,-11.2875)--(9.6,-12.899999999999999) node[lab,midway,below]{Ne pas publier};
\draw[br] (7.4,-9.675)--(9.6,-8.6) node[lab,midway,above]{Poursuivre};
\draw[br] (7.4,-9.675)--(9.6,-10.75) node[lab,midway,below]{Laisser faire};
\fill[nyt] (0.4,-7.928125) circle (4pt); \node[font=\small\bfseries,text=nyt,above left=3pt and -2pt] at (0.4,-7.928125) {Times};
\fill[gov] (2.6,-4.56875) circle (4pt); \node[font=\small\bfseries,text=gov,above left=3pt and -2pt] at (2.6,-4.56875) {Nixon};
\fill[wp] (5.0,-2.6875) circle (4pt); \node[font=\small\bfseries,text=wp,above left=3pt and -2pt] at (5.0,-2.6875) {Post};
\fill[gov] (7.4,-1.075) circle (4pt); \node[font=\small\bfseries,text=gov,above left=3pt and -2pt] at (7.4,-1.075) {Nixon};
\fill[wp] (5.0,-11.2875) circle (4pt); \node[font=\small\bfseries,text=wp,above left=3pt and -2pt] at (5.0,-11.2875) {Post};
\fill[gov] (7.4,-9.675) circle (4pt); \node[font=\small\bfseries,text=gov,above left=3pt and -2pt] at (7.4,-9.675) {Nixon};
\fill (9.6,0.0) circle (2pt);
\node[anchor=west,draw=black!35,fill=white,rounded corners=3pt,text width=15.2cm,align=left,font=\footnotesize,inner sep=4pt] at (9.85,0.0) {\textbf{$O_1$ · les deux journaux poursuivis}\\\textcolor{nyt}{\textbf{Times} : $p\,(B_T-c_T)+(1-p)\,q\,(B_T-c_T)+(1-p)(1-q)\,(\theta B_T-c_T-S_T)$}\\\textcolor{wp}{\textbf{Post} : $p\cdot0+(1-p)\,q\,(B_W-c_W)+(1-p)(1-q)\,(\theta B_W-c_W-S_W)$}\\\textcolor{gov}{\textbf{Nixon} : $p\,(-D-k_T)+(1-p)\,q\,(-D-k_T-k_W)+(1-p)(1-q)\,(-\theta D-k_T-k_W)$}};
\fill (9.6,-2.15) circle (2pt);
\node[anchor=west,draw=black!35,fill=white,rounded corners=3pt,text width=15.2cm,align=left,font=\footnotesize,inner sep=4pt] at (9.85,-2.15) {\textbf{$O_2$ · Times poursuivi, Post libre}\\\textcolor{nyt}{\textbf{Times} : $B_T-c_T$}\\\textcolor{wp}{\textbf{Post} : $p\cdot0+(1-p)\,B_W$}\\\textcolor{gov}{\textbf{Nixon} : $-D-k_T$}};
\fill (9.6,-4.3) circle (2pt);
\node[anchor=west,draw=black!35,fill=white,rounded corners=3pt,text width=15.2cm,align=left,font=\footnotesize,inner sep=4pt] at (9.85,-4.3) {\textbf{$O_3$ · Times poursuivi, Post muet}\\\textcolor{nyt}{\textbf{Times} : $p\,(B_T-c_T)+(1-p)\,(\theta B_T-c_T-S_T)$}\\\textcolor{wp}{\textbf{Post} : $0$}\\\textcolor{gov}{\textbf{Nixon} : $p\,(-D-k_T)+(1-p)\,(-\theta D-k_T)$}};
\fill (9.6,-6.449999999999999) circle (2pt);
\node[anchor=west,draw=black!35,fill=white,rounded corners=3pt,text width=15.2cm,align=left,font=\footnotesize,inner sep=4pt] at (9.85,-6.449999999999999) {\textbf{$O_4$ · Times publie, pas de procès}\\\textcolor{nyt}{\textbf{Times} : $B_T$}\\\textcolor{wp}{\textbf{Post} : $0$}\\\textcolor{gov}{\textbf{Nixon} : $-D$}};
\fill (9.6,-8.6) circle (2pt);
\node[anchor=west,draw=black!35,fill=white,rounded corners=3pt,text width=15.2cm,align=left,font=\footnotesize,inner sep=4pt] at (9.85,-8.6) {\textbf{$O_5$ · Post seul, poursuivi}\\\textcolor{nyt}{\textbf{Times} : $0$}\\\textcolor{wp}{\textbf{Post} : $p\,(B_W-c_W)+(1-p)\,(\theta B_W-c_W-S_W)$}\\\textcolor{gov}{\textbf{Nixon} : $p\,(-D-k_W)+(1-p)\,(-\theta D-k_W)$}};
\fill (9.6,-10.75) circle (2pt);
\node[anchor=west,draw=black!35,fill=white,rounded corners=3pt,text width=15.2cm,align=left,font=\footnotesize,inner sep=4pt] at (9.85,-10.75) {\textbf{$O_6$ · Post seul, pas de procès}\\\textcolor{nyt}{\textbf{Times} : $0$}\\\textcolor{wp}{\textbf{Post} : $B_W$}\\\textcolor{gov}{\textbf{Nixon} : $-D$}};
\fill (9.6,-12.899999999999999) circle (2pt);
\node[anchor=west,draw=black!35,fill=white,rounded corners=3pt,text width=15.2cm,align=left,font=\footnotesize,inner sep=4pt] at (9.85,-12.899999999999999) {\textbf{$O_7$ · personne ne publie}\\\textcolor{nyt}{\textbf{Times} : $0$}\\\textcolor{wp}{\textbf{Post} : $0$}\\\textcolor{gov}{\textbf{Nixon} : $0$}};
\end{tikzpicture}}
\end{landscape}


\subsection{Les deux règles}
\begin{align}
\text{Journal } j :\quad u_j&=\begin{cases}0 & \text{s'il ne publie pas,}\\ B_j\,\ind{ça sort}-c_j\,\ind{poursuivi}-S_j\,\ind{censuré} & \text{s'il publie,}\end{cases}\\[4pt]
\text{Nixon :}\quad u_N&=-D\,\ind{ça sort}-k_T\,\ind{Times poursuivi}-k_W\,\ind{Post poursuivi}.
\end{align}
Sans $\theta$, une censure veut dire que \textbf{rien ne sort} : le journal censuré ne touche aucun bénéfice et Nixon n'a pas de dommage.

\subsection{Les gains après un verdict}
\begin{center}
\small
\begin{tabular}{@{}lll@{\qquad}lll@{}}
\toprule
\multicolumn{3}{@{}l}{Journaux} & \multicolumn{3}{l@{}}{Nixon} \\
\midrule
\T{Times gagne} & $B_T-c_T$ & $9$ & Tout sort, 1 procès & $-D-k$ & $-201$ \\
\T{Times perd} & $-c_T-S_T$ & $-6$ & Tout sort, 2 procès & $-D-k_T-k_W$ & $-202$ \\
\W{Post gagne} & $B_W-c_W$ & $5$ & Censure, 1 procès & $-k$ & $-1$ \\
\W{Post perd} & $-c_W-S_W$ & $-10$ & Censure, 2 procès & $-k_T-k_W$ & $-2$ \\
\bottomrule
\end{tabular}
\end{center}

\subsection{Les sept feuilles}
\paragraph{$O_1$ — les deux journaux poursuivis.}
\begin{align*}
\T{\E[u_T]} &= p\,(B_T-c_T)+(1-p)\,q\,(B_T-c_T)+(1-p)(1-q)\,(-c_T-S_T)\\
&= 0,3\times9+0,455\times9+0,245\times(-6)=2,7+4,095-1,47=\mathbf{5,325}\\[4pt]
\W{\E[u_W]} &= p\times0+(1-p)\,q\,(B_W-c_W)+(1-p)(1-q)\,(-c_W-S_W)\\
&= 0+0,455\times5+0,245\times(-10)=2,275-2,45=\mathbf{-0,175}\\[4pt]
\N{\E[u_N]} &= p\,(-D-k_T)+(1-p)\,q\,(-D-k_T-k_W)+(1-p)(1-q)\,(-k_T-k_W)\\
&= 0,3\times(-201)+0,455\times(-202)+0,245\times(-2)=-60,3-91,91-0,49=\mathbf{-152,7}
\end{align*}

\paragraph{$O_2$ — Times poursuivi, Post libre.}
\begin{align*}
\T{\E[u_T]} &= B_T-c_T=\mathbf{9}\\
\W{\E[u_W]} &= p\times0+(1-p)\,B_W=0,7\times10=\mathbf{7}\\
\N{\E[u_N]} &= -D-k_T=\mathbf{-201}
\end{align*}

\paragraph{$O_3$ — Times poursuivi, Post muet.}
\begin{align*}
\T{\E[u_T]} &= p\,(B_T-c_T)+(1-p)(-c_T-S_T)=0,3\times9+0,7\times(-6)=2,7-4,2=\mathbf{-1,5}\\
\W{\E[u_W]} &= \mathbf{0}\\
\N{\E[u_N]} &= p\,(-D-k_T)+(1-p)(-k_T)=0,3\times(-201)+0,7\times(-1)=-60,3-0,7=\mathbf{-61}
\end{align*}

\paragraph{$O_4$, $O_6$, $O_7$ — pas de procès.}
\begin{align*}
O_4 &= (\T{B_T}\,;\,\W{0}\,;\,\N{-D})=(10\,;\,0\,;\,-200)\\
O_6 &= (\T{0}\,;\,\W{B_W}\,;\,\N{-D})=(0\,;\,10\,;\,-200)\\
O_7 &= (0\,;\,0\,;\,0)
\end{align*}

\paragraph{$O_5$ — Post seul, poursuivi.}
\begin{align*}
\T{\E[u_T]} &= \mathbf{0}\\
\W{\E[u_W]} &= p\,(B_W-c_W)+(1-p)(-c_W-S_W)=0,3\times5+0,7\times(-10)=1,5-7=\mathbf{-5,5}\\
\N{\E[u_N]} &= p\,(-D-k_W)+(1-p)(-k_W)=\mathbf{-61}
\end{align*}

\begin{landscape}
\begin{center}{\large\bfseries L'arbre corrigé, en chiffres, sans $\theta$}\\[2pt]{\small Gains dans l'ordre \textcolor{nyt}{\textbf{Times}}, \textcolor{wp}{\textbf{Post}}, \textcolor{gov}{\textbf{Nixon}}. $p$ = juge local, $q$ = Cour suprême. Exemple : $p=0,3$, $q=0,65$. Traits rouges épais : choix à l'équilibre $\Rightarrow$ $O_7$, rien ne sort.}\end{center}\vspace{-4pt}
\noindent\resizebox{\linewidth}{!}{\begin{tikzpicture}[x=1cm,y=1cm,lab/.style={font=\scriptsize,fill=white,inner sep=1pt,sloped},eq/.style={line width=1.8pt,draw=gov},br/.style={line width=0.6pt,draw=black!55}]
\draw[br] (0.4,-7.928125)--(2.6,-4.56875) node[lab,midway,above]{Publier};
\draw[eq] (0.4,-7.928125)--(5.0,-11.2875) node[lab,midway,below]{Ne pas publier};
\draw[eq] (2.6,-4.56875)--(5.0,-2.6875) node[lab,midway,above]{Poursuivre};
\draw[br] (2.6,-4.56875)--(9.6,-6.449999999999999) node[lab,midway,below]{Laisser faire};
\draw[br] (5.0,-2.6875)--(7.4,-1.075) node[lab,midway,above]{Publier};
\draw[eq] (5.0,-2.6875)--(9.6,-4.3) node[lab,midway,below]{Ne pas publier};
\draw[eq] (7.4,-1.075)--(9.6,0.0) node[lab,midway,above]{Poursuivre};
\draw[br] (7.4,-1.075)--(9.6,-2.15) node[lab,midway,below]{Laisser faire};
\draw[br] (5.0,-11.2875)--(7.4,-9.675) node[lab,midway,above]{Publier};
\draw[eq] (5.0,-11.2875)--(9.6,-12.899999999999999) node[lab,midway,below]{Ne pas publier};
\draw[eq] (7.4,-9.675)--(9.6,-8.6) node[lab,midway,above]{Poursuivre};
\draw[br] (7.4,-9.675)--(9.6,-10.75) node[lab,midway,below]{Laisser faire};
\fill[nyt] (0.4,-7.928125) circle (4pt); \node[font=\small\bfseries,text=nyt,above left=3pt and -2pt] at (0.4,-7.928125) {Times};
\fill[gov] (2.6,-4.56875) circle (4pt); \node[font=\small\bfseries,text=gov,above left=3pt and -2pt] at (2.6,-4.56875) {Nixon};
\fill[wp] (5.0,-2.6875) circle (4pt); \node[font=\small\bfseries,text=wp,above left=3pt and -2pt] at (5.0,-2.6875) {Post};
\fill[gov] (7.4,-1.075) circle (4pt); \node[font=\small\bfseries,text=gov,above left=3pt and -2pt] at (7.4,-1.075) {Nixon};
\fill[wp] (5.0,-11.2875) circle (4pt); \node[font=\small\bfseries,text=wp,above left=3pt and -2pt] at (5.0,-11.2875) {Post};
\fill[gov] (7.4,-9.675) circle (4pt); \node[font=\small\bfseries,text=gov,above left=3pt and -2pt] at (7.4,-9.675) {Nixon};
\fill (9.6,0.0) circle (2pt);
\node[anchor=west,draw=black!35,fill=white,rounded corners=3pt,text width=15.2cm,align=left,font=\footnotesize,inner sep=4pt] at (9.85,0.0) {\textbf{$O_1$ · les deux journaux poursuivis}\\\textcolor{nyt}{\textbf{Times} : $p\cdot9+(1-p)\,q\cdot9+(1-p)(1-q)\cdot(-6)$ $=\mathbf{5,325}$}\\\textcolor{wp}{\textbf{Post} : $p\cdot0+(1-p)\,q\cdot5+(1-p)(1-q)\cdot(-10)$ $=\mathbf{-0,175}$}\\\textcolor{gov}{\textbf{Nixon} : $p\cdot(-201)+(1-p)\,q\cdot(-202)+(1-p)(1-q)\cdot(-2)$ $=\mathbf{-152,7}$}};
\fill (9.6,-2.15) circle (2pt);
\node[anchor=west,draw=black!35,fill=white,rounded corners=3pt,text width=15.2cm,align=left,font=\footnotesize,inner sep=4pt] at (9.85,-2.15) {\textbf{$O_2$ · Times poursuivi, Post libre}\\\textcolor{nyt}{\textbf{Times} : $9$}\\\textcolor{wp}{\textbf{Post} : $p\cdot0+(1-p)\cdot10$ $=\mathbf{7}$}\\\textcolor{gov}{\textbf{Nixon} : $-201$}};
\fill (9.6,-4.3) circle (2pt);
\node[anchor=west,draw=black!35,fill=white,rounded corners=3pt,text width=15.2cm,align=left,font=\footnotesize,inner sep=4pt] at (9.85,-4.3) {\textbf{$O_3$ · Times poursuivi, Post muet}\\\textcolor{nyt}{\textbf{Times} : $p\cdot9+(1-p)\cdot(-6)$ $=\mathbf{-1,5}$}\\\textcolor{wp}{\textbf{Post} : $0$}\\\textcolor{gov}{\textbf{Nixon} : $p\cdot(-201)+(1-p)\cdot(-1)$ $=\mathbf{-61}$}};
\fill (9.6,-6.449999999999999) circle (2pt);
\node[anchor=west,draw=black!35,fill=white,rounded corners=3pt,text width=15.2cm,align=left,font=\footnotesize,inner sep=4pt] at (9.85,-6.449999999999999) {\textbf{$O_4$ · Times publie, pas de procès}\\\textcolor{nyt}{\textbf{Times} : $10$}\\\textcolor{wp}{\textbf{Post} : $0$}\\\textcolor{gov}{\textbf{Nixon} : $-200$}};
\fill (9.6,-8.6) circle (2pt);
\node[anchor=west,draw=black!35,fill=white,rounded corners=3pt,text width=15.2cm,align=left,font=\footnotesize,inner sep=4pt] at (9.85,-8.6) {\textbf{$O_5$ · Post seul, poursuivi}\\\textcolor{nyt}{\textbf{Times} : $0$}\\\textcolor{wp}{\textbf{Post} : $p\cdot5+(1-p)\cdot(-10)$ $=\mathbf{-5,5}$}\\\textcolor{gov}{\textbf{Nixon} : $p\cdot(-201)+(1-p)\cdot(-1)$ $=\mathbf{-61}$}};
\fill (9.6,-10.75) circle (2pt);
\node[anchor=west,draw=black!35,fill=white,rounded corners=3pt,text width=15.2cm,align=left,font=\footnotesize,inner sep=4pt] at (9.85,-10.75) {\textbf{$O_6$ · Post seul, pas de procès}\\\textcolor{nyt}{\textbf{Times} : $0$}\\\textcolor{wp}{\textbf{Post} : $10$}\\\textcolor{gov}{\textbf{Nixon} : $-200$}};
\fill (9.6,-12.899999999999999) circle (2pt);
\node[anchor=west,draw=black!35,fill=white,rounded corners=3pt,text width=15.2cm,align=left,font=\footnotesize,inner sep=4pt] at (9.85,-12.899999999999999) {\textbf{$O_7$ · personne ne publie}\\\textcolor{nyt}{\textbf{Times} : $0$}\\\textcolor{wp}{\textbf{Post} : $0$}\\\textcolor{gov}{\textbf{Nixon} : $0$}};
\end{tikzpicture}}
\end{landscape}
\subsection{La résolution, en partant de la fin}
\begin{center}
\small
\begin{tabular}{@{}lll@{}}
\toprule
Qui décide & Comparaison & Choix \\
\midrule
\N{Nixon} face au Post en relais & $O_1:-152,7$ contre $O_2:-201$ & poursuit \\
\W{Post} en relais & $O_1:-0,175$ contre $O_3:0$ & \textbf{se tait} \\
\N{Nixon} face au Times & $O_3:-61$ contre $O_4:-200$ & poursuit \\
\N{Nixon} face au Post seul & $O_5:-61$ contre $O_6:-200$ & poursuit \\
\W{Post} seul & $O_5:-5,5$ contre $O_7:0$ & se tait \\
\T{Times} & $O_3:-1,5$ contre $O_7:0$ & \textbf{ne publie pas} \\
\bottomrule
\end{tabular}
\end{center}
\begin{resultat}
\textbf{Résultat sans $\theta$ : $O_7$, rien n'est publié.} Ce n'est pas l'histoire de 1971.
\end{resultat}

\subsection{Pourquoi : les seuils}
\textbf{Le Post en relais} publie si son gain espéré est positif. On divise par $1-p>0$ :
\[
q\,(B_W-c_W)+(1-q)(-c_W-S_W)\ge0\iff q\ \ge\ \frac{c_W+S_W}{B_W+S_W}=\frac{5+5}{10+5}=\frac23\approx66,7\,\%.
\]
Avec $q=65\,\%$, il est juste en dessous : il se tait.

\textbf{Le Times}, resté seul face au juge local, publie si
\[
p\,(B_T-c_T)+(1-p)(-c_T-S_T)\ge0\iff p\ \ge\ \frac{c_T+S_T}{B_T+S_T}=\frac{1+5}{10+5}=40\,\%.
\]
Avec $p=30\,\%$, il renonce.

\textbf{Nixon} poursuit presque toujours : un procès lui coûte $1$ et peut lui éviter $200$. Par exemple, il poursuit le Times resté seul si $p<1-\dfrac{k_T}{D}=99,5\,\%$.

\begin{simple}
Sans $\theta$, perdre son procès efface \emph{tout} le bénéfice. Les journaux risquent alors beaucoup pour un gain de seulement $10$ : il leur faut de très bonnes chances de gagner.
\end{simple}

\section{Avec \texorpdfstring{$\theta$}{theta} : la part déjà publiée avant la censure}

\subsection{Ce qui change dans les équations}
Une censure n'arrête la publication qu'après les premiers articles. En 1971, le Times avait publié trois articles avant d'être stoppé. Seules les lignes « censure » changent :
\begin{align*}
\text{Journal censuré :}\quad & \theta B_j-c_j-S_j \quad\text{au lieu de}\quad -c_j-S_j,\\
\text{Nixon après une censure :}\quad & -\theta D-k \quad\text{au lieu de}\quad -k .
\end{align*}
Les règles deviennent, avec $\sigma=1$ (tout sort), $\theta$ (censure) ou $0$ (rien) :
\begin{align*}
u_j &= B_j\,\sigma-c_j\,\ind{poursuivi}-S_j\,\ind{censuré}\quad(\text{s'il publie ; }0\text{ sinon}),\\
u_N &= -D\,\sigma-k_T\,\ind{Times poursuivi}-k_W\,\ind{Post poursuivi}.
\end{align*}

\subsection{Les gains après une censure, avec \texorpdfstring{$\theta=0,3$}{theta = 0,3}}
\begin{center}
\small
\begin{tabular}{@{}lll@{}}
\toprule
Situation & Calcul & Valeur \\
\midrule
\T{Times perd} & $0,3\times10-1-5$ & $-3$ \quad (au lieu de $-6$) \\
\W{Post perd} & $0,3\times10-5-5$ & $-7$ \quad (au lieu de $-10$) \\
\N{Nixon}, censure, 1 procès & $-0,3\times200-1$ & $-61$ \quad (au lieu de $-1$) \\
\N{Nixon}, censure, 2 procès & $-0,3\times200-2$ & $-62$ \quad (au lieu de $-2$) \\
\bottomrule
\end{tabular}
\end{center}

\subsection{Les feuilles qui changent}
\begin{align*}
O_1:\quad \T{\E[u_T]} &= p(B_T-c_T)+(1-p)q(B_T-c_T)+(1-p)(1-q)(\theta B_T-c_T-S_T)\\
&= 2,7+4,095+0,245\times(-3)=\mathbf{6,06}\\
\W{\E[u_W]} &= (1-p)q(B_W-c_W)+(1-p)(1-q)(\theta B_W-c_W-S_W)=2,275+0,245\times(-7)=\mathbf{0,56}\\
\N{\E[u_N]} &= p(-D-k_T)+(1-p)q(-D-k_T-k_W)+(1-p)(1-q)(-\theta D-k_T-k_W)\\
&= -60,3-91,91+0,245\times(-62)=\mathbf{-167,4}\\[6pt]
O_3:\quad \T{\E[u_T]} &= p(B_T-c_T)+(1-p)(\theta B_T-c_T-S_T)=2,7+0,7\times(-3)=\mathbf{0,6}\\
\N{\E[u_N]} &= p(-D-k_T)+(1-p)(-\theta D-k_T)=-60,3+0,7\times(-61)=\mathbf{-103}\\[6pt]
O_5:\quad \W{\E[u_W]} &= p(B_W-c_W)+(1-p)(\theta B_W-c_W-S_W)\\ &=1,5+0,7\times(-7)=\mathbf{-3,4}\\
\N{\E[u_N]} &= p(-D-k_W)+(1-p)(-\theta D-k_W)=\mathbf{-103}
\end{align*}
$O_2$, $O_4$, $O_6$ et $O_7$ ne changent pas : aucune censure n'y tient.

\begin{landscape}
\begin{center}{\large\bfseries L'arbre corrigé, en chiffres, avec $\theta=0,3$}\\[2pt]{\small Gains dans l'ordre \textcolor{nyt}{\textbf{Times}}, \textcolor{wp}{\textbf{Post}}, \textcolor{gov}{\textbf{Nixon}}. $p$ = juge local, $q$ = Cour suprême. Exemple : $p=0,3$, $q=0,65$. Traits rouges épais : choix à l'équilibre $\Rightarrow$ $O_1$, c'est 1971.}\end{center}\vspace{-4pt}
\noindent\resizebox{\linewidth}{!}{\begin{tikzpicture}[x=1cm,y=1cm,lab/.style={font=\scriptsize,fill=white,inner sep=1pt,sloped},eq/.style={line width=1.8pt,draw=gov},br/.style={line width=0.6pt,draw=black!55}]
\draw[eq] (0.4,-7.928125)--(2.6,-4.56875) node[lab,midway,above]{Publier};
\draw[br] (0.4,-7.928125)--(5.0,-11.2875) node[lab,midway,below]{Ne pas publier};
\draw[eq] (2.6,-4.56875)--(5.0,-2.6875) node[lab,midway,above]{Poursuivre};
\draw[br] (2.6,-4.56875)--(9.6,-6.449999999999999) node[lab,midway,below]{Laisser faire};
\draw[eq] (5.0,-2.6875)--(7.4,-1.075) node[lab,midway,above]{Publier};
\draw[br] (5.0,-2.6875)--(9.6,-4.3) node[lab,midway,below]{Ne pas publier};
\draw[eq] (7.4,-1.075)--(9.6,0.0) node[lab,midway,above]{Poursuivre};
\draw[br] (7.4,-1.075)--(9.6,-2.15) node[lab,midway,below]{Laisser faire};
\draw[br] (5.0,-11.2875)--(7.4,-9.675) node[lab,midway,above]{Publier};
\draw[eq] (5.0,-11.2875)--(9.6,-12.899999999999999) node[lab,midway,below]{Ne pas publier};
\draw[eq] (7.4,-9.675)--(9.6,-8.6) node[lab,midway,above]{Poursuivre};
\draw[br] (7.4,-9.675)--(9.6,-10.75) node[lab,midway,below]{Laisser faire};
\fill[nyt] (0.4,-7.928125) circle (4pt); \node[font=\small\bfseries,text=nyt,above left=3pt and -2pt] at (0.4,-7.928125) {Times};
\fill[gov] (2.6,-4.56875) circle (4pt); \node[font=\small\bfseries,text=gov,above left=3pt and -2pt] at (2.6,-4.56875) {Nixon};
\fill[wp] (5.0,-2.6875) circle (4pt); \node[font=\small\bfseries,text=wp,above left=3pt and -2pt] at (5.0,-2.6875) {Post};
\fill[gov] (7.4,-1.075) circle (4pt); \node[font=\small\bfseries,text=gov,above left=3pt and -2pt] at (7.4,-1.075) {Nixon};
\fill[wp] (5.0,-11.2875) circle (4pt); \node[font=\small\bfseries,text=wp,above left=3pt and -2pt] at (5.0,-11.2875) {Post};
\fill[gov] (7.4,-9.675) circle (4pt); \node[font=\small\bfseries,text=gov,above left=3pt and -2pt] at (7.4,-9.675) {Nixon};
\fill (9.6,0.0) circle (2pt);
\node[anchor=west,draw=black!35,fill=white,rounded corners=3pt,text width=15.2cm,align=left,font=\footnotesize,inner sep=4pt] at (9.85,0.0) {\textbf{$O_1$ · les deux journaux poursuivis}\\\textcolor{nyt}{\textbf{Times} : $p\cdot9+(1-p)\,q\cdot9+(1-p)(1-q)\cdot(-3)$ $=\mathbf{6,06}$}\\\textcolor{wp}{\textbf{Post} : $p\cdot0+(1-p)\,q\cdot5+(1-p)(1-q)\cdot(-7)$ $=\mathbf{0,56}$}\\\textcolor{gov}{\textbf{Nixon} : $p\cdot(-201)+(1-p)\,q\cdot(-202)+(1-p)(1-q)\cdot(-62)$ $=\mathbf{-167,4}$}};
\fill (9.6,-2.15) circle (2pt);
\node[anchor=west,draw=black!35,fill=white,rounded corners=3pt,text width=15.2cm,align=left,font=\footnotesize,inner sep=4pt] at (9.85,-2.15) {\textbf{$O_2$ · Times poursuivi, Post libre}\\\textcolor{nyt}{\textbf{Times} : $9$}\\\textcolor{wp}{\textbf{Post} : $p\cdot0+(1-p)\cdot10$ $=\mathbf{7}$}\\\textcolor{gov}{\textbf{Nixon} : $-201$}};
\fill (9.6,-4.3) circle (2pt);
\node[anchor=west,draw=black!35,fill=white,rounded corners=3pt,text width=15.2cm,align=left,font=\footnotesize,inner sep=4pt] at (9.85,-4.3) {\textbf{$O_3$ · Times poursuivi, Post muet}\\\textcolor{nyt}{\textbf{Times} : $p\cdot9+(1-p)\cdot(-3)$ $=\mathbf{0,6}$}\\\textcolor{wp}{\textbf{Post} : $0$}\\\textcolor{gov}{\textbf{Nixon} : $p\cdot(-201)+(1-p)\cdot(-61)$ $=\mathbf{-103}$}};
\fill (9.6,-6.449999999999999) circle (2pt);
\node[anchor=west,draw=black!35,fill=white,rounded corners=3pt,text width=15.2cm,align=left,font=\footnotesize,inner sep=4pt] at (9.85,-6.449999999999999) {\textbf{$O_4$ · Times publie, pas de procès}\\\textcolor{nyt}{\textbf{Times} : $10$}\\\textcolor{wp}{\textbf{Post} : $0$}\\\textcolor{gov}{\textbf{Nixon} : $-200$}};
\fill (9.6,-8.6) circle (2pt);
\node[anchor=west,draw=black!35,fill=white,rounded corners=3pt,text width=15.2cm,align=left,font=\footnotesize,inner sep=4pt] at (9.85,-8.6) {\textbf{$O_5$ · Post seul, poursuivi}\\\textcolor{nyt}{\textbf{Times} : $0$}\\\textcolor{wp}{\textbf{Post} : $p\cdot5+(1-p)\cdot(-7)$ $=\mathbf{-3,4}$}\\\textcolor{gov}{\textbf{Nixon} : $p\cdot(-201)+(1-p)\cdot(-61)$ $=\mathbf{-103}$}};
\fill (9.6,-10.75) circle (2pt);
\node[anchor=west,draw=black!35,fill=white,rounded corners=3pt,text width=15.2cm,align=left,font=\footnotesize,inner sep=4pt] at (9.85,-10.75) {\textbf{$O_6$ · Post seul, pas de procès}\\\textcolor{nyt}{\textbf{Times} : $0$}\\\textcolor{wp}{\textbf{Post} : $10$}\\\textcolor{gov}{\textbf{Nixon} : $-200$}};
\fill (9.6,-12.899999999999999) circle (2pt);
\node[anchor=west,draw=black!35,fill=white,rounded corners=3pt,text width=15.2cm,align=left,font=\footnotesize,inner sep=4pt] at (9.85,-12.899999999999999) {\textbf{$O_7$ · personne ne publie}\\\textcolor{nyt}{\textbf{Times} : $0$}\\\textcolor{wp}{\textbf{Post} : $0$}\\\textcolor{gov}{\textbf{Nixon} : $0$}};
\end{tikzpicture}}
\end{landscape}
\subsection{La résolution}
\begin{center}
\small
\begin{tabular}{@{}lll@{}}
\toprule
Qui décide & Comparaison & Choix \\
\midrule
\N{Nixon} face au Post en relais & $O_1:-167,4$ contre $O_2:-201$ & poursuit \\
\W{Post} en relais & $O_1:0,56$ contre $O_3:0$ & \textbf{publie} \\
\N{Nixon} face au Times & $O_1:-167,4$ contre $O_4:-200$ & poursuit \\
\N{Nixon} face au Post seul & $O_5:-103$ contre $O_6:-200$ & poursuit \\
\W{Post} seul & $O_5:-3,4$ contre $O_7:0$ & se tait \\
\T{Times} & $O_1:6,06$ contre $O_7:0$ & \textbf{publie} \\
\bottomrule
\end{tabular}
\end{center}
\begin{resultat}
\textbf{Résultat avec $\theta=0,3$ : $O_1$.} Le Times publie, Nixon le poursuit, le Post prend le relais, Nixon le poursuit aussi : c'est 1971.
\end{resultat}

\subsection{Les seuils avec \texorpdfstring{$\theta$}{theta}}
\begin{align*}
\text{Post en relais :}\quad & q\ \ge\ \frac{c_W+S_W-\theta B_W}{(1-\theta)B_W+S_W}=\frac{5+5-3}{7+5}=\frac{7}{12}\approx58,3\,\%\quad(\text{au lieu de }66,7\,\%),\\[4pt]
\text{Times :}\quad & \pi\ \ge\ \frac{c_T+S_T-\theta B_T}{(1-\theta)B_T+S_T}=\frac{1+5-3}{7+5}=25\,\%\quad(\text{au lieu de }40\,\%),\\[4pt]
\text{Nixon contre le Post :}\quad & q\ <\ 1-\frac{k_W}{(1-\theta)D}=1-\frac{1}{140}\approx99,3\,\%,\\[4pt]
\text{Nixon contre le Times :}\quad & k_T+(1-p)\,k_W\ <\ (1-p)(1-q)(1-\theta)D,\quad\text{soit}\quad 1,7<34,3 .
\end{align*}
Ici $\pi=p+(1-p)q=0,755$ est la chance que le Times gagne à la fin quand le Post prend le relais.

\section{Comparaison et recommandation}
\begin{center}
\small
\begin{tabular}{@{}lcc@{}}
\toprule
 & Sans $\theta$ & Avec $\theta=0,3$ \\
\midrule
$O_1$ (Times ; Post ; Nixon) & $(5,325\ ;\ -0,175\ ;\ -152,7)$ & $(6,06\ ;\ 0,56\ ;\ -167,4)$ \\
$O_3$ & $(-1,5\ ;\ 0\ ;\ -61)$ & $(0,6\ ;\ 0\ ;\ -103)$ \\
$O_5$ & $(0\ ;\ -5,5\ ;\ -61)$ & $(0\ ;\ -3,4\ ;\ -103)$ \\
Seuil du Post en relais & $q\ge66,7\,\%$ & $q\ge58,3\,\%$ \\
Seuil du Times & $40\,\%$ & $25\,\%$ \\
Résultat & $O_7$ : rien ne sort & $O_1$ : 1971 \\
\bottomrule
\end{tabular}
\end{center}
\begin{simple}
\textbf{Recommandation : garder $\theta$.} C'est historiquement juste (trois articles parus avant la censure), et sans lui une censure efface tout, ce qui rend la publication trop risquée pour un bénéfice de $10$.
\end{simple}

\section{Les points à trancher ensemble}
\begin{enumerate}[leftmargin=*,itemsep=3pt]
\item \textbf{Le coût du Post} : $3$ (dessin) ou $5$ (liste) ? Avec $3$ et sans $\theta$, le seuil du Post devient $\frac{3+5}{10+5}\approx53,3\,\%$ et il prendrait le relais.
\item \textbf{Le dommage de Nixon} : $5$ (dessin) ou $200$ (liste) ? C'est décisif. Avec $D=5$ et $\theta=0,3$, Nixon ne poursuit plus le Times : le coût moyen des procès ($1,7$) dépasse le dommage qu'il évite ($0,7\times0,35\times0,7\times5\approx0,86$). Tout sort sans procès ($O_4$). Pour retrouver 1971, $D$ doit être grand par rapport au coût des procès.
\item \textbf{La sanction $S=5$} : la garder ? La même pour les deux journaux ? Dans notre modèle principal, le Post risque plus à cause de son entrée en bourse et de ses licences de télévision.
\item \textbf{$O_2$} : le Times censuré garde-t-il $9$ si le Post a tout publié ? (Ici : oui.)
\item \textbf{$O_1$} : le Post a-t-il $0$ si le juge local libère le Times ? (Ici : oui.)
\item \textbf{Les probabilités} $p$ et $q$ ne sont pas fixées. Sans $\theta$, il faut $q\ge66,7\,\%$ pour que le Post prenne le relais.
\end{enumerate}

\end{document}
